% Pursuing Stacks (cotes 134-2 to 134-8): archive cote 134-8, site page 25 (PDF page 26), typescript page 579.
% typescript. Draft compared: ps7.tex:1157-1213.
% AI-transcribed by Claude Opus 5.5 (claude-opus-5-5, adaptive thinking, effort high), Message Batches API, batch ps4 (msgbatch_012Aoc1NuskATPsyBeKSYK7J), 2026-11-12; not yet checked by eye.
\scanpage{134-8/25}
\mspage{579}

is reduced to a point  -  visibly not a very "cellular" behaviour when
$\hand{|}I_1\hand{|}$ is'nt aspheric !  At any rate,  as the sheaves \xout{$R^if_*$ may be compu}
$R^i\hand{|}f\hand{|}_*(F)$  (for F a constant sheaf above say) may be computed fiberwise,
we see that if $I_1$  is'nt aspheric,  these sheaves (which are constant
on $\hand{|}J\hand{|}\hand{\setminus\{}1\hand{\}}$ )  are \uline{not} going to be constant on  $\hand{|}J\hand{|}\hand{\setminus\{}0\hand{\}}$, as \xout{we expect}
\xout{for} they should if we want an algebraic paradigm of operations like
$R\hand{|}f\hand{|}_*$  in terms of sheaves on finite ordered sets.

In this example we could take $\hand{|}I\hand{|}$ to be a perfect n-cell (n$\hand{\geq}$1),
hence  $I_1$  is an (n-1)-sphere,  whereas $\hand{|}J\hand{|}$ is'nt really a (combinatori-%
al) 1-cell,  as it's boundary has just \uline{one} point 0 ,  instead of two.
The 1-cell structure corresponds to the ordered set (\uline{opposite} to the
ordered set formed by the two endpoints and the dimension 1 stratum)
\diagram{\[ J = \begin{tikzcd}[row sep=small, column sep=small] & \hand{y} \ar[dl, "\hand{}"'] \ar[dr, "\hand{}"] & \\ 0 & & 1 \end{tikzcd} \quad . \]}
For any ordered set I,  to give a map  f: I $\hand{\longrightarrow}$ J  amount to give
the two subsets
\[ I_0 = f^{-1}(0) \ , \ \ \deleted{\xout{\unsure{et}}}\ I_1 = f^{-1}(1) \qquad , \]
subject\xout{s} to the only condition of being open,  and disjoint.  In terms
of strata of $\hand{|}I\hand{|}$ ,  this means that we give \over{to}{two} sets $I_0,I_1$ of strata,
containing with any stratum any smaller one,  and \over{with}{having} no stratum \inserted{in}
common \deleted{to $I_0$ and $I_1$}.  The condition that the sets (13) should be non-%
empty says that any point in I which is neither in $I_0$ nor $I_1$ , admits
majorants in both  -  or geometrically,  any stratum which is'nt in
$I_0$ nor in $I_1$ meets both $\hand{|}I_0\hand{|}$ and $\hand{|}I_1\hand{|}$ i.e. contains strata which are
in $I_0$ and strata which are in $I_1$ .   Even when $\hand{|}I\hand{|}$ is a combinatorial
2-cell i.e. a polygonal disc,  this condition does not imply asphérici-%
ty of the sets (13) (not even 0-connectedness).  To see this,  we take
a polygon of even order 2m $\hand{\geq}$\over{\illegible{1}}{4} ,  and we take  $I_0$ and $I_1$ to consist
of the vertices of the polygon of even resp. odd index.  We take the
set (13) where x is the dimension 2 stratum  i.e. the \xout{mini} smallest
element of I ,  mapped to y (the smallest in J),  and y' either 0 or 1.
The reader will easily figure out on a drawing the structure of the
map $\hand{|}f\hand{|}$ ,  as I just did myself:  the fibers at the endpoints of the
segment $\hand{|}J\hand{|}$ are,  as given,  discrete with cardinal m,  the fiber at
a point different from the endpoints and \inserted{from} their barycenter are disjoint
sums of m segments (hence homotopic to the former fibers),  whereas
the fiber at the barycenter is the union of m segments meeting in their

\margin{25}
